\documentclass{article}
\usepackage{amsmath, amssymb, geometry, hyperref}
\geometry{margin=1in}
\setlength{\parskip}{0.8em}

\begin{document}


\section*{Overview}
To bring coins from Coinbase to a local wallet, you are \textbf{withdrawing cryptocurrency from Coinbase} and sending it to your own \textbf{wallet address}. The exact steps depend on the type of local wallet you are using.

\section{Local Software Wallet (e.g., Electrum, Exodus, MetaMask)}
\subsection*{1. Open Your Local Wallet}
\begin{itemize}
    \item Ensure the wallet is installed and fully synced (if required).
    \item Locate the \textbf{Receive} section.
    \item Copy your \textbf{wallet address} (a long alphanumeric string).
\end{itemize}
\textbf{Note:} Always ensure the coin type matches. For example, send BTC only to a Bitcoin address, not to an Ethereum address.

\subsection*{2. Access Coinbase}
\begin{itemize}
    \item Log in to \href{https://www.coinbase.com}{Coinbase}.
    \item On the dashboard, click \textbf{``Send \& Receive''}.
    \item Under the \textbf{Send} tab, choose the cryptocurrency you want to withdraw.
\end{itemize}

\subsection*{3. Enter Your Wallet Address}
\begin{itemize}
    \item Paste your local wallet address into the ``To'' field.
    \item Enter the amount you want to send.
    \item Double-check the address and coin type.
\end{itemize}

\subsection*{4. Choose Network (If Applicable)}
For tokens that exist on multiple networks (e.g., USDT on ERC20, TRC20, or Base):
\begin{itemize}
    \item Select the same network that your local wallet supports.
    \item Coinbase will display a warning if the network is incompatible.
\end{itemize}

\subsection*{5. Confirm and Send}
\begin{itemize}
    \item Review all transaction details carefully.
    \item Confirm the send using your 2FA or other security method.
    \item Coinbase will broadcast the transaction to the blockchain.
\end{itemize}

\subsection*{6. Wait for Confirmations}
\begin{itemize}
    \item Transfers typically take from a few minutes to an hour.
    \item You can track progress using the blockchain explorer link provided by Coinbase.
\end{itemize}

\section{Local Hardware Wallet (e.g., Ledger, Trezor)}
\begin{itemize}
    \item Use Ledger Live or Trezor Suite to generate your receive address.
    \item Plug in and unlock your device before initiating the transfer.
    \item Verify the address directly on the device screen before confirming on Coinbase.
\end{itemize}

\section{Example: Sending ETH from Coinbase to MetaMask}
\begin{enumerate}
    \item Open MetaMask, click ``Account 1,'' and copy your Ethereum address (\texttt{0x...}).
    \item On Coinbase, click ``Send \& Receive'' and select Ethereum.
    \item Paste your MetaMask address, enter the amount, and send.
    \item Wait a few minutes and check MetaMask after confirmation.
\end{enumerate}

\section{Important Tips}
\begin{itemize}
    \item Always test with a small amount first.
    \item Never send coins between incompatible networks (e.g., Base $\rightarrow$ Ethereum without bridging).
    \item Be aware of network fees, especially on Ethereum.
    \item Consider using \textbf{Coinbase Wallet} (self-custodial) if you are new to managing private keys.
\end{itemize}


\newpage


\section*{Overview}
We compare three things:
\begin{itemize}
  \item Ledger (hardware wallet)  
  \item MetaMask (software wallet)  
  \item A normal SSD (storage device)  
\end{itemize}
We also include how Ledger supports XRP (Ripple) as described on the Ledger XRP wallet page.

\section{What They Are}

\subsection*{Ledger (Hardware Wallet)}
\begin{itemize}
  \item A physical device (often USB, potentially Bluetooth) that stores **private keys** in secure hardware offline.  
  \item It signs transactions internally, so your keys never leave the device.  
  \item According to the Ledger XRP wallet page, Ledger supports XRP by allowing users to add an XRP account in the Ledger Live app and manage XRP (send, receive, swap) while the private keys remain securely stored in hardware. \\
  \(\quad\) It states: “Ledger hardware wallet stores your private keys and signs transactions offline, making them resistant to malicious attacks.” :contentReference[oaicite:0]{index=0} \\
  \(\quad\) Also, Ledger Live (the software companion) acts as the interface to manage XRP and other assets. :contentReference[oaicite:1]{index=1}  
  \item The XRP wallet feature is marketed by Ledger as “secure your XRP” with hardware-backed security. :contentReference[oaicite:2]{index=2}  
  \item The Ledger page notes that XRP addresses begin with the letter “r” and are case-sensitive. :contentReference[oaicite:3]{index=3}  
\end{itemize}

\subsection*{MetaMask (Software Wallet)}
\begin{itemize}
  \item A software application (browser extension or mobile) which stores your private key (or seed phrase) locally (often encrypted) on your device.  
  \item It enables you to interact with Ethereum, EVM-compatible chains, and dApps.  
  \item Because it is “hot” (connected or accessible by the device), it has greater exposure to software attacks compared to a hardware wallet.  
\end{itemize}

\subsection*{Normal SSD (Storage Device)}
\begin{itemize}
  \item A solid-state drive is hardware for general data storage (files, programs, operating systems).  
  \item It is not designed for cryptographic key security by itself.  
  \item If you store private keys (or wallet files) on an SSD, they are only as secure as your system’s protections (encryption, backups, malware defense).  
  \item The SSD has no native concept of “signing transactions” or interfacing with blockchains — it is simply storage.
\end{itemize}

\section{Key Differences and Roles}

\begin{tabular}{|l|l|l|l|}
\hline
Feature & Ledger (Hardware) & MetaMask (Software) & SSD (Storage) \\
\hline
Purpose & Secure offline key storage \& transaction signing & Wallet interface and key management & General file/data storage \\
\hline
Private keys stored & In secure hardware, isolated & Encrypted on device storage & Not inherently — if stored here, vulnerable \\
\hline
Connectivity & Offline (cold), only connected when needed & Online (hot) & Always connected when in use \\
\hline
Risk & Low (physical access + PIN needed) & Moderate (software attacks, device compromise) & High (if used to store keys) \\
\hline
Transaction signing & Done internally in device & Done by software on the device & Not applicable \\
\hline
Asset storage & Keys only (actual coins live on blockchain) & Keys only & Raw storage (no blockchain interface) \\
\hline
Supports XRP? & Yes (via Ledger + Ledger Live) & Only if wallet software supports XRP protocol & N/A \\
\hline
\end{tabular}

\section{How Ledger Supports XRP (Ripple)}

From the Ledger XRP wallet page:

\begin{itemize}
  \item Ledger offers an XRP wallet solution (hardware + Ledger Live) to store and manage XRP securely. :contentReference[oaicite:4]{index=4}  
  \item Users can add an XRP account in Ledger Live and use it to send, receive, swap XRP, while private keys stay in hardware. :contentReference[oaicite:5]{index=5}  
  \item XRP addresses are case-sensitive and begin with “r”. :contentReference[oaicite:6]{index=6}  
  \item Ledger’s site emphasizes that hardware wallets keep private keys off connected systems, protecting them from many online threats. :contentReference[oaicite:7]{index=7}  
\end{itemize}

\section{Conclusion}

\begin{itemize}
  \item **Ledger** is a hardware wallet that gives you strong security by keeping keys offline, and Ledger supports XRP (Ripple) via its integrated wallet solution.  
  \item **MetaMask** is a software (hot) wallet, convenient but more exposed.  
  \item A **normal SSD** is a generic data storage device, not built with key security or blockchain functionality in mind.  
\end{itemize}




\end{document}
